% Abstract for: Finding Bugs in Python Native Extensions via Safety Obligation Checking
\newcommand{\paperabstract}{%
Python native extensions written in C/C++ are widely used to provide computational acceleration, hardware interaction, and system-level functionality for performance-critical Python libraries. Developing such extensions requires complying with the safety requirements of the Python/C API, including reference ownership management, null-value handling, exception propagation, and resource lifecycle management. However, these requirements are dispersed throughout extensive CPython documentation and are difficult to reason about across complex execution paths, causing bugs such as memory leaks, use-after-free vulnerabilities, and interpreter crashes that often elude conventional static analysis. We make the key observation that Python/C API safety requirements can be systematically modeled as \emph{dischargeable safety obligations}.
Based on a systematic study of CPython documentation and hundreds of real-world bugs, we derive 13 obligation categories that capture the dominant classes of Python/C API misuse and formally define the conditions under which each obligation is discharged.
We propose \tool, a bug detection framework for Python native extensions based on safety-obligation checking. \tool first employs lightweight static analysis to identify obligation-generating API usages and eliminate candidates lacking evidence of violations.
It then leverages an LLM to perform path-sensitive reasoning over obligation creation and discharge, determining whether the corresponding discharge conditions hold on every execution path. We evaluate \tool on 9 widely used open-source projects. \tool finds 86 previously unknown bugs with 82.7\% precision, of which 45 have been confirmed or fixed by maintainers.%
}
